\begin{abstract}
A Latin square is FFF if every two-line permutation in its row,
column and symbol views has only even cycles. We establish structural
constraints and construction results for this property. Direct products
take the componentwise Boolean OR of the three failure patterns and are
group-isotopic precisely when both factors are. Subsquare and congruence
obstructions restrict constructions, while explicit tables of orders
$12$ and $14$ disprove the former power-of-two existence conjecture.
The order-$10$ exclusion combines a group-theoretic reduction with exact
finite computation. Finite-field, row-dependent and block-substitution
constructions yield further infinite families. A fixed finite audit
distinguishes $1703$ FFF main classes of order $20$; affine-orbit counting
then gives lower bounds in higher orders, not censuses. A pattern-preserving
corner trade produces examples with exactly one intercalate. For orders
congruent to $2$ modulo $4$, a triangle of near-involutive line pairs
forces an odd cycle in a companion view. In a binary-fibre extension,
transversals project to simple two-plexes with an exact parity criterion
and lift count; disjoint transversals with the same projection cannot
produce an FFF double prolongation. Written proofs, finite computational
dependencies and bounded search results are distinguished. The unrestricted
existence question at order $18$, the complete order spectrum and the
literature priority of the recent constructions remain open.
\end{abstract}
\noindent\textbf{Keywords.} Latin square; quasigroup; Latin trade; direct
product; main class; affine orbit; computational enumeration.
\par\smallskip
\noindent\textbf{2020 Mathematics Subject Classification.}
Primary 05B15; Secondary 20N05, 05C70.
